\chapter{Introduction}
\label{chapterlabel1}

\section{Research Background}

In recent years, the ``24-hour city'' has become an increasingly important concept in urban policy. London is a prominent example: the introduction of the Night Tube in 2016, together with subsequent policies on the night-time economy, culture and governance \citep{TfL2016,GLAEconomics2024}, signals that city authorities no longer treat the night simply as a pause in daytime activity, but increasingly recognise it as a period with independent social and economic value \citep{Shaw2010,Shaw2014}.

Against this backdrop, around one quarter of London's workforce is estimated to work during night-time hours \citep{Mavrogeni2025}, and the mobility demands generated by night-time employment, consumption and leisure have become an important challenge for the transport system.

However, urban transport planning has long been organised around daytime activity and commuting, shaping service design, infrastructure allocation and the evaluation frameworks used to assess transport systems \citep{McArthur2019}. Academic research on how cities function after dark has also remained relatively limited \citep{VanLiempt2015}. This daytime orientation in both planning and research is particularly problematic because nighttime travel is not simply a temporal extension of daytime travel. The transport environment at night differs: public transport services are typically more limited, lighting and safety conditions change, and travellers' purposes and resources are more diverse \citep{Kapitza2022,Kapitza2025}. Night-time transport must serve not only leisure and social activities, but also the mobility of night-shift workers and other time-constrained groups \citep{Plyushteva2019,Shaw2022,Palm2023}. These characteristics create particular challenges for understanding when and where night-time transport demand occurs.

High-temporal-resolution smart-card data and clustering methods have been widely used to identify typical usage patterns in public transport systems \citep{Briand2016,ElMahrsi2017,Manley2018,Gan2020,Zhao2020}. However, these analyses have typically focused on weekday commuting, peak-hour flows or all-day ridership structures. The night-time period has more often appeared as part of the full daily cycle than as a distinct analytical window.

Thus, it remains unclear whether stable and recurring temporal usage patterns exist within night-time public transport, and how these patterns are distributed across urban space.

This study addresses this gap by analysing high-temporal-resolution smart-card data for London's rail and bus systems. It identifies temporal patterns of night-time public transport use, examines their spatial distribution, and situates them against the geography of London's night-time workforce \citep{Mavrogeni2025} and broader socio-spatial context.

\section{Research Question}

Building on this background, the study poses the following research question:

Using high-temporal-resolution public transport usage data, what distinct patterns of night-time public transport usage can we observe within London's rail and bus systems, and what are their temporal and spatial characteristics? What differences exist in the night-time usage patterns of these two modes of transport, and how do they relate to the geographical distribution of night-time work in London and the broader socio-spatial context at the area level?

This question is addressed through two specific sub-questions:

\textbf{RQ1}: What distinct night-time usage types can be identified for rail stations and bus-served LSOAs, based on their temporal, directional and day type activity patterns? How are these types distributed across London and in what respects do the two transport modes differ?

\textbf{RQ2}: To what extent are the identified usage types associated with the London Night Workers Classification (LNWC) and with local-level socio-economic, household, housing and employment contexts? How do these associations help us understand the urban environments under the different usage types?

To address these questions, the study first identifies night-time usage types separately within the Rail and Bus systems, then characterises their temporal, directional, day-type and spatial patterns, and finally examines how the fixed usage types are associated with LNWC and wider area-level urban contexts. This sequence moves from the identification of observed transport-use patterns to their external contextual interpretation.

This study contributes by relating night-time usage types derived from high-temporal-resolution transport data to an independently constructed geography of night-time work and the broader socio-spatial context. This provides both a methodological link between temporal transport analysis and area-based contextual data, and an empirical account of how night-time transport patterns are situated across London.

\section{Dissertation Structure}

The remainder of the dissertation is organised as follows. Chapter 2 reviews literature on the urban night, night-time mobility and smart-card clustering methods. Chapter 3 describes the study area, data, feature construction, clustering methodology, and contextual analysis. Chapter 4 presents the identified rail and bus usage types, their temporal and spatial characteristics, and their associations with the LNWC and wider socio-spatial conditions. Chapter 5 interprets these findings in relation to the research questions and wider literature, considers their implications for night-time transport research and planning, and reflects on the main methodological limitations. Chapter 6 concludes the dissertation by synthesising the answers to the research questions and summarising the study's main contribution.
